% =====================================================================
% 02-Spectral.tex
% Geometric Hallucination Taxonomy — Mathematical Spine
% Borealis S. Hedling — 2026
% =====================================================================

\section*{02. Spectral Operator}

This module formalizes the spectral derivative, eigenvalue dynamics, and
instability criteria used to define Spectral Cascade Hallucinations (SCH)
in the Geometric Hallucination Taxonomy (GHT). The definitions rely on
the geometric foundations established in 00-Foundations.tex.

% ---------------------------------------------------------------------
\subsection*{2.1 Latent Representation Matrices}

Let $A_t$ denote the latent representation at inference step $t$. We
assume $A_t$ is a real, square matrix of size $n \times n$:
\[
A_t \in \mathbb{R}^{n \times n}.
\]

The matrix $A_t$ may arise from:
\begin{itemize}
\item Jacobians of update rules,
\item local linearizations of latent dynamics,
\item attention or mixing operators,
\item synthetic constructions in controlled experiments.
\end{itemize}

The spectral operator analyzes the eigenvalues of $A_t$.

% ---------------------------------------------------------------------
\subsection*{2.2 Eigenvalues}

Let $\lambda_1(t), \dots, \lambda_n(t)$ be the eigenvalues of $A_t$,
counted with algebraic multiplicity:
\[
A_t v_i(t) = \lambda_i(t) v_i(t).
\]

We assume $A_t$ is diagonalizable for simplicity; the instability
criteria extend to the general case via Jordan normal form.

% ---------------------------------------------------------------------
\subsection*{2.3 Continuous Spectral Derivative}

For a smooth family of matrices $A(t)$, the continuous spectral
derivative is
\[
\Delta_\lambda(t)
= \frac{d}{dt} \lambda(t).
\]

This derivative measures instantaneous spectral drift.

% ---------------------------------------------------------------------
\subsection*{2.4 Discrete Spectral Derivative}

In discrete inference settings, we define
\[
\Delta_\lambda(t)
= \lambda(t+1) - \lambda(t).
\]

This is the quantity used in the Python toolkit and experimental
modules.

% ---------------------------------------------------------------------
\subsection*{2.5 Spectral Norm Drift}

The spectral norm of $A_t$ is
\[
\|A_t\|_2 = \max_i |\lambda_i(t)|.
\]

Spectral norm drift is defined as
\[
\Delta_{\|A\|}(t)
= \|A_{t+1}\|_2 - \|A_t\|_2.
\]

Large positive drift indicates explosive growth; large negative drift
indicates collapse.

% ---------------------------------------------------------------------
\subsection*{2.6 Spectral Cascade Instability}

A spectral cascade occurs when eigenvalues exhibit runaway behavior.

We define the instability criterion:
\[
|\Delta_\lambda(t)| > \sigma_{\mathrm{thr}}.
\]

A cascade is present if this condition holds for infinitely many $t$.

Explosive cascade:
\[
\lambda(t) \to +\infty.
\]

Collapse cascade:
\[
\lambda(t) \to 0.
\]

Oscillatory cascade:
\[
\lambda(t+1) \cdot \lambda(t) < 0.
\]

These behaviors correspond to distinct SCH subtypes.

% ---------------------------------------------------------------------
\subsection*{2.7 Spectral Stability Region}

Define the spectral stability region
\[
\Omega_{\sigma}
= \{ A \in \mathbb{R}^{n \times n} :
|\lambda_i(A_{t+1}) - \lambda_i(A_t)| \le \sigma_{\mathrm{thr}}
\text{ for all } i \}.
\]

A latent trajectory is spectrally stable if $A_t \in \Omega_{\sigma}$
for all $t$.

% ---------------------------------------------------------------------
\subsection*{2.8 Spectral Cascades and Latent Geometry}

Spectral cascades correspond to geometric instability in the latent
manifold. If $A_t$ represents a local linearization of the update rule
$\Phi$, then eigenvalue drift indicates instability in the local
geometry of the trajectory.

In particular:
\begin{itemize}
\item explosive eigenvalues correspond to local expansion,
\item collapsing eigenvalues correspond to local contraction,
\item oscillatory eigenvalues correspond to alternating curvature regimes.
\end{itemize}

These geometric interpretations motivate the SCH classification.

% ---------------------------------------------------------------------
\subsection*{2.9 Discrete Approximation}

In numerical settings, eigenvalues are computed via standard algorithms
(e.g., QR decomposition). The discrete spectral derivative is computed
directly from successive eigenvalue lists.

The Python toolkit implements this approximation and validates it in
\texttt{/tests/test\_spectral.py}.

% ---------------------------------------------------------------------
\subsection*{2.10 Spectral–Cascade Hallucinations (SCH)}

A latent trajectory exhibits SCH if there exists an eigenvalue sequence
$\lambda(t)$ such that
\[
|\Delta_\lambda(t)| > \sigma_{\mathrm{thr}}
\quad \text{for infinitely many } t.
\]

Equivalently, SCH occurs when spectral instability prevents the latent
trajectory from remaining within a stable linearization regime.

% =====================================================================
% End of 02-Spectral.tex
% =====================================================================
